\begin{mistake}{2026-06-11}{02}

\begin{problemcontent}
设 \(f''(x) < 0\)，\(f(1) = 2\)，\(f'(1) = -3\)，求证：\(f(x) = 0\) 在 \((1, +\infty)\) 有且仅有一个实根。
\end{problemcontent}

\begin{solutioncontent}
\textbf{唯一性}：
因为 \(f''(x) < 0\) 对任意 \(x \in (1, +\infty)\) 均成立，所以 \(f'(x)\) 在 \((1, +\infty)\) 上严格单调递减。
又因为 \(f'(1) = -3 < 0\)，所以对于任意 \(x > 1\)，都有：
\[ f'(x) < f'(1) = -3 < 0 \]
因此 \(f(x)\) 在 \((1, +\infty)\) 上严格单调递减。
单调函数与 \(x\) 轴最多只有一个交点，故 \(f(x) = 0\) 在 \((1, +\infty)\) 内最多只有一个实根。

\textbf{存在性}（提供两种方法）：
\textbf{方法一（泰勒公式）}：
将 \(f(x)\) 在 \(x_0 = 1\) 处展开至一阶（带有拉格朗日余项的泰勒公式）：
\[ f(x) = f(1) + f'(1)(x-1) + \frac{f''(\xi)}{2}(x-1)^2 \]
其中 \(1 < \xi < x\)。代入已知条件得：
\[ f(x) = 2 - 3(x-1) + \frac{f''(\xi)}{2}(x-1)^2 \]
因为 \(f''(x) < 0\)，所以二阶项 \(\frac{f''(\xi)}{2}(x-1)^2 < 0\)，放缩得到：
\[ f(x) < 2 - 3(x-1) \]
取 \(x = 2\)，则 \(f(2) < 2 - 3(2-1) = -1 < 0\)。
又因为 \(f(1) = 2 > 0\) 且 \(f(x)\) 连续，由零点定理可知，在区间 \((1, 2)\) 内至少存在一点 \(c\) 使得 \(f(c) = 0\)。

\textbf{方法二（拉格朗日中值定理）}：
在区间 \([1, x]\) 上对 \(f(t)\) 使用拉格朗日中值定理：
\[ f(x) - f(1) = f'(\eta)(x-1) \]
其中 \(\eta \in (1, x)\)。因为 \(f'(x)\) 严格单调递减，所以 \(f'(\eta) < f'(1) = -3\)。
由于 \(x > 1\)，两边同乘 \((x-1)\) 得：
\[ f'(\eta)(x-1) < -3(x-1) \]
即 \(f(x) - 2 < -3(x-1)\)，得到放缩：
\[ f(x) < 2 - 3(x-1) \]
同样取 \(x = 2\) 得到 \(f(2) < 0\)。结合 \(f(1) = 2 > 0\)，由零点定理得证在 \((1,2)\) 内至少存在一个实根。

综合可知，\(f(x) = 0\) 在 \((1, +\infty)\) 有且仅有一个实根。
\end{solutioncontent}

\mistakeknowledge{
\begin{itemize}
  \item \textbf{核心公式/定理}：零点定理（存在性）、单调性（唯一性）、带有拉格朗日余项的泰勒公式、拉格朗日中值定理。
  \item \textbf{易错点}：证明宏观区间的存在性或不等式必须使用带有拉格朗日余项的泰勒公式，绝不能使用皮亚诺余项。
  \item \textbf{解题技巧}：已知函数值、导数值及二阶导符号时，利用泰勒公式或中值定理将函数值放缩为其切线（一次函数）来寻找异号点。
\end{itemize}}

\end{mistake}
